% ----------------------------------------------------------------------
\subsubsection*{ کاربرد \lr{Clean Architecture} در طراحی این پروژه}

ایده‌ی اصلی \lr{Clean Architecture}: (۱) قوانین کسب‌وکار باید از جزئیات فناوری جدا باشند، و (۲) جهت وابستگی‌ها فقط رو به داخل، یعنی به‌سمت \lr{Entities} و \lr{Use Cases} است. هیچ کلاسی در لایه‌های داخلی حق ندارد به یک کتابخانه یا فریمورک خارجی اشاره کند.

\begin{itemize}
    \item \textbf{لایه‌ی \lr{Entities (Domain)}:} موجودیت \lr{Message}، \lr{Value Object}ای به نام \lr{QueueStatus}، \lr{port}های \lr{Queue}، \lr{QueueManager}، \lr{QueueFactory} و \lr{Logger}، و مجموعه‌ای از خطاهای کسب‌وکار به‌صورت \lr{sentinel error} (\lr{ErrQueueFull}، \lr{ErrQueueEmpty}، \lr{ErrQueueNotFound}، \lr{ErrQueueAlreadyExists}، ...). این لایه \emph{هیچ} \lr{import} خارجی ندارد.
    \item \textbf{لایه‌ی \lr{UseCases}:} کلاس \lr{QueueUseCase} منطق برنامه را مدیریت می‌کند و دو \lr{port} خروجی به نام‌های \lr{IDGenerator} و \lr{Clock} تعریف می‌کند تا حتی به \lr{uuid} یا \lr{time.Now} \emph{مستقیم} وابسته نباشد. این لایه از طریق \lr{port} ورودی \lr{QueueService} توسط لایه‌ی \lr{Delivery} مصرف می‌شود.
    \item \textbf{لایه‌ی \lr{Interface Adapters} - بخش \lr{Repository}:} پیاده‌سازی \lr{in-memory} صف به‌صورت \lr{Ring Buffer} (\lr{RingBufferQueue}) و یک \lr{QueueManager} امن از نظر هم‌زمانی (با \lr{RWMutex}). \lr{QueueManager} از طریق یک \lr{QueueFactory} تزریق‌شده، نوع صف زیرین را نمی‌داند؛ یعنی می‌توان بعدا \lr{Redis-backed Queue} اضافه کرد بدون اینکه این کلاس تغییر کند.
    \item \textbf{لایه‌ی \lr{Interface Adapters} - بخش \lr{Delivery}:} \lr{HTTP} \lr{handler}های مبتنی بر \lr{Gin}، \lr{DTO}ها، یک \lr{router} مستقل، و یک تابع متمرکز \lr{writeDomainError} که خطاهای دامنه را به کد وضعیت \lr{HTTP} نگاشت می‌کند. تمام این لایه تنها به \lr{interface} \lr{usecase.QueueService} وابسته است، نه به نوع concrete .
    \item \textbf{لایه‌ی \lr{Service}:} \lr{TTLCleaner} که با یک \lr{ticker} و \lr{context.Context} پیام‌های منقضی را پاکسازی می‌کند و قابلیت \lr{graceful shutdown} دارد. این \lr{worker} فقط به پورت‌های \lr{QueueManager} و \lr{Logger} وابسته است.
    \item \textbf{لایه‌ی \lr{Frameworks \& Drivers}:} \lr{idgen.UUID}، \lr{clock.Real}، \lr{logger.Zap} و خود \lr{Gin}. این‌ها فقط در \lr{Composition Root} (یعنی \lr{cmd/server/main.go}) به لایه‌های داخلی تزریق می‌شوند.
\end{itemize}

\textbf{رعایت اصل \lr{DIP}:} لایه‌ی \lr{UseCase} \lr{port} موردنیاز خود را \emph{خودش} تعریف می‌کند (\lr{IDGenerator}, \lr{Clock}) و لایه‌ی بیرونی موظف است آن را پیاده کند. 

نتیجه: تعویض \lr{uuid} با هر تولیدکننده‌ی \lr{ID} دیگر، یا تعویض \lr{Gin} با \lr{Echo/Fiber/gRPC}، تنها یک \lr{adapter} جدید می‌خواهد و دست‌کاری در منطق کسب‌وکار لازم نیست.


\textbf{نکات کلیدی پیاده‌سازی:}
\begin{itemize}
    \item پیاده‌سازی صف به‌صورت \lr{Ring Buffer} با پیچیدگی $O(1)$ برای \lr{push} و \lr{pop} و بدون \lr{allocation} اضافه به ازای هر پیام.
    \item انقضای پیام به دو شیوه‌ی موازی: (۱) \lr{Lazy} هنگام \lr{Pop}: پیام‌های منقضی‌شده در ابتدای صف رد می‌شوند، (۲) \lr{Proactive} از طریق \lr{TTLCleaner} که هر یک ثانیه یک‌بار اجرا می‌شود تا حتی پیام‌هایی که هرگز \lr{pop} نمی‌شوند نیز حافظه را نگه ندارند.
    \item قرارداد: \lr{TTL=0} به معنی پیام بدون تاریخ انقضاست.
    \item همه‌ی خطاهای کسب‌وکار به‌صورت مقدارهای \lr{sentinel} (\lr{errors.New}) در پکیج \lr{domain} تعریف شده‌اند و فقط در یک نقطه‌ی متمرکز (تابع \lr{writeDomainError}) به کد وضعیت \lr{HTTP} نگاشت می‌شوند. این یعنی هیچ مفهوم \lr{HTTP}ای داخل لایه‌های داخلی نشت نمی‌کند.
    \item \lr{Composition Root} در \lr{cmd/server/main.go} سیگنال‌های \lr{SIGINT/SIGTERM} را با \lr{signal.NotifyContext} می‌گیرد، سپس \lr{TTLCleaner} و \lr{http.Server} را با مهلت ۵ ثانیه‌ای \lr{Shutdown} می‌کند.
\end{itemize}

\textbf{مسیرهای \lr{HTTP} پیاده‌سازی‌شده:}

\begin{latin}
\begin{verbatim}
GET    /healthz                # liveness probe
POST   /queues                 # create queue (capacity required)
GET    /queues                 # list queues
GET    /queues/:queue          # status: {id, name, size, capacity}
POST   /queues/:queue/push     # enqueue {value, ttl}
POST   /queues/:queue/pop      # dequeue; 204 if empty
\end{verbatim}
\end{latin}

\textbf{کتاب‌خانه‌ها و دلیل انتخاب آن‌ها:}
\begin{itemize}
    \item \lr{github.com/gin-gonic/gin}: \lr{HTTP router} سبک با \lr{middleware} منعطف. فقط در لایه‌ی \lr{delivery} استفاده می‌شود.
    \item \lr{github.com/google/uuid}: تولید \lr{UUID v4} برای شناسه‌ی صف‌ها و پیام‌ها. در پشت \lr{port} \lr{IDGenerator} ایزوله شده.
    \item \lr{go.uber.org/zap}: \lr{logger} ساخت‌یافته و سریع. در پشت \lr{port} \lr{domain.Logger} ایزوله شده و فقط در \lr{internal/pkg/logger} حضور دارد.
\end{itemize}

% ----------------------------------------------------------------------

\textbf{اجرای پروژه:}
\begin{latin}
\begin{verbatim}
cd Baby-Rabbit
go run ./cmd/server          # server listens on :8080
\end{verbatim}
\end{latin}


\textbf{تست با \lr{Postman}:} فایل \lr{Baby-Rabbit.postman\_collection.json} شامل تمام \lr{request}ها به‌همراه \lr{test script}های \lr{JavaScript} است که به‌صورت خودکار شناسه‌ی صف ساخته‌شده را در \lr{collection variable} ذخیره می‌کند. کافی است \lr{collection} را \lr{import} کنید و درخواست‌ها را به ترتیب اجرا کنید.

\textbf{اعتبارسنجی قاعده‌ی وابستگی:} می‌توان به‌صورت تجربی نشان داد که هیچ‌یک از لایه‌های داخلی به فریمورک‌های بیرونی وابسته نیستند:

\begin{latin}
\begin{verbatim}
go list -f '{{join .Imports "\n"}}' ./internal/domain
# (empty — pure Go)

go list -f '{{join .Imports "\n"}}' ./internal/usecase
# Baby-Rabbit/internal/domain
# strings
# time

go list -f '{{join .Imports "\n"}}' ./internal/service
# Baby-Rabbit/internal/domain
# context
# time
\end{verbatim}
\end{latin}

به‌عبارت دیگر، با گرفتن یک \lr{snapshot} از \lr{import graph} می‌بینیم که هیچ ارجاعی از \lr{domain/usecase/service} به \lr{gin}، \lr{uuid}، \lr{zap} یا حتی \lr{net/http} وجود ندارد.

% ----------------------------------------------------------------------
\subsubsection*{ تحلیل نهایی معماری}

اگر معماری به‌دست‌آمده را با دوایر هم‌مرکز اصلی معماری \lr{Clean} مقایسه کنیم، تطابق به‌صورت زیر است:

\begin{itemize}
    \item \lr{Entities} $\Leftrightarrow$ بسته‌ی \lr{domain} (\lr{Message} و \lr{port}های \lr{Queue}/\lr{QueueManager}).
    \item \lr{Use Cases} $\Leftrightarrow$ بسته‌ی \lr{usecase} (\lr{QueueUseCase} و \lr{port}های آن).
    \item \lr{Interface Adapters} $\Leftrightarrow$ بسته‌های \lr{delivery/http} و \lr{repository}.
    \item \lr{Frameworks \& Drivers} $\Leftrightarrow$ \lr{gin}، \lr{uuid}، \lr{zap}، \lr{net/http} که فقط در لایه‌ی بیرونی و در \lr{Composition Root} ظاهر می‌شوند.
\end{itemize}

\textbf{قاعده‌ی وابستگی برقرار است:} هیچ فایلی در \lr{domain} و \lr{usecase} از پکیج‌های \lr{gin}، \lr{uuid} یا \lr{zap} \lr{import} نمی‌کند و این موضوع را با \lr{go list -deps} می‌توان به‌صورت آبجکتیو تأیید کرد.

\textbf{نزدیک‌ترین الگو:} از منظر الگوها، معماری نهایی به \lr{Hexagonal Architecture} (که با نام \lr{Ports \& Adapters} نیز شناخته می‌شود) بسیار نزدیک است. دلیل این نزدیکی روشن است: \lr{port}های ورودی (\lr{QueueService}) و \lr{port}های خروجی (\lr{IDGenerator}، \lr{Clock}، \lr{QueueManager}، \lr{Queue}، \lr{QueueFactory}، \lr{Logger}) به‌صورت صریح در پکیج‌های داخلی تعریف شده‌اند و هر \lr{adapter} (\lr{Gin HTTP}، \lr{Ring Buffer}، \lr{UUID}، \lr{System Clock}، \lr{Zap}) دقیقاً روی لبه‌ی شش‌ضلعی قرار دارد. در عمل، \lr{Clean Architecture} و \lr{Hexagonal} در همین پروژه به یک‌جا می‌رسند: تفاوت آن‌ها بیشتر در ادبیات (\lr{layers} در برابر \lr{hexagon}) است تا در ساختار نهایی کد. اگر بخواهیم دقیق‌تر بگوییم، چون ما در پایین مرز شش‌ضلعی هیچ \lr{driving adapter}ای (مثل \lr{CLI} یا \lr{message consumer}) نداشتیم و تنها یک \lr{HTTP adapter} ورودی داریم، حالت ساده‌شده‌ی \lr{Hexagonal} است.

\textbf{مزایای ملموسی که این طراحی به ما داد:}
\begin{itemize}
    \item جایگزینی فناوری در لایه‌ی بیرونی (مثلاً مهاجرت از \lr{Gin} به \lr{Echo} یا حتی به \lr{gRPC}) تنها همان لایه را دست‌خوش تغییر می‌کند.
    \item جایگزینی \lr{Ring Buffer} با هر \lr{backing store} دیگری (مثلاً \lr{Redis} یا \lr{BadgerDB}) فقط یک \lr{QueueFactory} تازه می‌خواهد و یک خط در \lr{main.go}. لایه‌های \lr{usecase} و \lr{domain} دست‌نخورده می‌مانند.
    \item جایگزینی \lr{zap} با \lr{slog} یا \lr{logr} فقط یک \lr{adapter} تازه برای \lr{domain.Logger} می‌خواهد.
    \item منطق برنامه بدون نیاز به \lr{HTTP server}، \lr{uuid} یا \lr{wall clock} قابل تست است و این موضوع در پوشش تست‌های واحد منعکس شده است.
\end{itemize}

